\documentclass[12pt]{article}

\usepackage[dvips,letterpaper,margin=0.75in,bottom=0.5in]{geometry}
\usepackage{cite}
\usepackage{slashed}
\usepackage{graphicx}
\usepackage{amsmath}
\usepackage{amssymb}
\usepackage{braket}
\usepackage{pythonhighlight}
\usepackage{mathtools}
\usepackage{tikz}
\begin{document}


\title{PHY 112L \\ A 2-D Ideal-Gas Atmosphere}
\author{Michael Mulhearn}

\maketitle

\section{Introduction}

In this assignment, we will model an atmosphere as an ideal gas in two
dimensions.  We will measure the heat exchanged during the approach to
thermal equilibrium, numerically determine the heat capacity of an ideal gas, and compare with our theoretical expectaion.  Next, we will add gravity to our model and measure how this impacts the molecular density and heat capacity of the atmosphere.

\section{Heat Capacity of an Ideal Gas}

For an ideal gas (without gravity) the total energy $U$ is purely
kinetic, and the equipartition theorem implies that:
\begin{equation}
U = \frac{f}{2}NkT
\end{equation}
where $f$ is the number of quadratic degrees of freedom: 3 for an
ideal-gas in 3-D, and 2 for an ideal-gas in 2-D.

The heat capacity $C$ describes quantitatively how the temperature of the gas changes as heat is added or removed:
\begin{equation}
C \equiv \frac{Q}{\Delta T}
\end{equation}
\noindent
Under constant volume, no work is done on the gas, and so:
$$\Delta U = Q + W = Q$$
Therefore, we can calculate the heat capacity under constant volume as:
$$C_V \equiv \left( \frac{Q}{\Delta T}\right)_V = \left( \frac{\Delta U}{\Delta T}\right)_V
$$
and taking $\Delta T$ small, we have:
\begin{equation}
C_V = \left( \frac{\partial U}{\partial T} \right)_V
\end{equation}
Applying this to an ideal gas we find that:
\begin{equation}
\label{eqn:cv}
C_V = \frac{f}{2} N k
\end{equation}

When two systems, $A$ and $B$, initally at temperatures $T_A$ and $T_B$ are brought into thermal contact, they reach an equilibrium temperature $T$ by exchanging heat.  The conservation of energy implies that:
$$C_A T_A + C_B T_B = C_A T + C_B T$$
where $C_A$ and $C_B$ are the heat capacities of the two sytems.  Equivalently:
\begin{equation}
T = \frac{C_A}{C_A + C_B} \cdot T_A + \frac{C_B}{C_A + C_B} \cdot T_B
\end{equation}
Notice that for an infinite reservoir $R$ with $C_R = \infty$ this reduces to:
$$T = T_R$$
Also, in the case that $C_A = C_B$, we have:
$$T = \frac{T_A + T_B}{2}$$
that is, the temperatue is just the average of the two system.

\section{Ideal Gas with Gravity}

Suppose now that our ideal gas is in a constant gravitation field with acceleration $g$.
Then each particle has gravitational potential energy:
\begin{equation}
E_{\rm grav} = m g y
\end{equation}
where we have set $y=0$ as the zero potential.  We expect the position
of the particles to follow a Boltzmann distribution, so that:
\begin{equation*}
p(y) \propto \exp\left( -\frac{E_{\rm grav}(y)}{kT} \right) = \exp\left( -\frac{m g y}{kT} \right)
\end{equation*}
The normalization condition:
$$\int_0^\infty p(y) dy = 1$$
gives us the normalization constant, and so:
\begin{equation}
p(y) = \frac{mg}{kT} \exp\left( -\frac{m g y}{kT} \right) 
\end{equation}
This is a decaying exponential, so that the atmosphere gets thinner as you go higher, exactly as we expect.  The average height of the gas molecules is:
\begin{equation}
\label{eqn:yexp}
\braket{y} = \frac{k T }{m g}
\end{equation}
The total gravitation potential energy for the entire gas is:
$$U_{\rm grav} = \sum_i m g y_i = N m g \braket{y}$$
and using Equation~\ref{eqn:yexp} we have:
$$U_{\rm grav} = N kT$$
The total energy $U$ of the system includes the kinetic energy part, so:
\begin{equation}
U = \frac{f}{2} N k T + NkT = \frac{f+2}{2} N k T
\end{equation}
thus the heat capacity of our atmosphere is:
\begin{equation}
\label{eqn:cvg}  
C_V = \left( \frac{\partial U}{\partial T} \right)_V = \frac{f+2}{2} N k 
\end{equation}
Notice the system is still at constant V (just V is now infinite) or,
more importantly, no external work is done to the atmosphere, so our
previous result for the kinetic part still apply.  The influence of
gravity increases the heat capacity by $Nk$.  This additional heat
capacity comes from the systems ability to stash additional energy as
gravitation potential energy.  This situation is quite analogous to a
gas under constant pressure, for which the same addtional factor for the heat capacity is obtained.

\section{System of Units}

We will use our same system of units as in the previous assignment
with one additional constraint.  Interesting, the ideal gas model of
HW 1 had no specific length scale defined, we simply took it to be
some scale of interest which we called $a$.  Our results apply equally
well to the air in child's balloon or the vast extremely low-density
plasmas of the interstellar medium.  Once we turn on gravity, however,
there is a relevant length scale given by Equation~\ref{eqn:yexp}.  So
we set:
\begin{equation}
 \label{eqn:newa}
a \equiv \frac{k T_0}{m g}
\end{equation}
so that:
\begin{equation}
  \label{eqn:expy}
  \braket{y/a} = \frac{T}{T_0}
\end{equation}
There is no need to change any of our results from the previous
homework.  {\bf The quantity $a$ will still be set to $1$ in your
  code.}  It's only in interpreting our results that we would be
forced to apply Equation~\ref{eqn:newa}.  

The other units that we defined in the previous homework continue to
apply, so let's recall them here.  Motivated by the Maxwell-Boltzmann
distribution, the velocity unit
$v_0$ is defined by:
$$ v_0^2 \equiv \frac{k T_0}{m}$$
and the unit for time is:
$$ \tau \equiv a / v_0$$
Notice that our units of acceleration are now:
$$\frac{a}{\tau^2} \equiv g$$
Motivated by the ideal gas law, we can set the unit of pressure to:
$$O_0 \equiv \frac{NkT_0}{a^2}$$
and the unit of total energy:
$$U_0 = N k T_0$$
Remember that for analytic work, we keep these factors in place.  But when using an
analytic result within our computer programs, we set all of these quantities to 1:
$$m = a = g = T_0 = v_0 = \tau = O_0 = U_0 = 1$$
However, if you would like, you can turn off gravity in a specific
context by setting $g=0$. Just take care that you turn gravity on with $g=1$.
As before, the temperature of an ideal gas is still calculated through
the average kinetic energy, which in our system of units becomes:
\begin{equation}
 \label{eqn:expv}
  \frac{T}{T_0} = \frac{\braket{v_x^2/v_0^2} + \braket{v_y^2/v_0^2}}{2}
\end{equation}
This is also the total kinetic energy of the gas, in units of total energy.
The total energy of the gas has a new term from the gravitation potentail energy, so that:
\begin{equation}
  \frac{U}{U_0} = \frac{\braket{v_x^2/v_0^2} + \braket{v_y^2/v_0^2}}{2} + \braket{y/a} 
\end{equation}  
Using Equations \ref{eqn:expy} and \ref{eqn:expv} we can write this rather elegantly as:
\begin{equation}
  \frac{U}{U_0} = 2 \cdot \frac{T}{T_0}
\end{equation}  
Recalling that:
$$C_V = \left( \frac{\partial U}{\partial T}\right)_V$$
we see that in our units, the heat capacity of our gas under the influence of gravity is:
\begin{equation}
\frac{C_V}{U_0 / T_0} = 2
\end{equation}
If we turn off gravity, we lose the gravitational potential term in the total energy and have:
\begin{equation}
  \frac{U}{U_0} = \frac{T}{T_0}
\end{equation}  
and:
\begin{equation}
\frac{C_V}{U_0 / T_0} = 1
\end{equation}
Notice that you can also get these results for $C_V$ from Equations~\ref{eqn:cv} and
\ref{eqn:cvg} if you simply divide by
$$U_0/T_0 = N k$$
and set $f=2$.

\section{Bouncing under Gravity}

When simulating our atmosphere under the influence of gravity, we'll
have to adjust our algorithm for circulating the gas and bouncing off
the ground.  We'll no longer need a top for our cylinder of gas, we'll
allow particles to reach as high as they can under the influence of
gravity.

Our numerical model includes the $y$-position of each molecules.
While inside the container, the molecules position can be updated for
a time step $\Delta t$ as:
\begin{equation}
y \to y + v_y \; \Delta t - \frac{1}{2} \; g \; (\Delta t)^2
\end{equation}
and the velocity by:
\begin{equation}
v_y \to v_y \; - g \Delta t
\end{equation}




\begin{figure}[htbp]
  \begin{center}
    
    \begin{tikzpicture}[scale=0.30]
      % Parameters
      \def\xa{5}    % x-intercept a
      \def\xb{35}   % x-intercept b
      \def\yv{5}     % vertex y-position

      % Compute k from x1, x2, yv: k = -4*yv / (x2-x1)^2
      \pgfmathsetmacro{\k}{-4*\yv/((\xb-\xa)^2)}

      % Parabola domain
      \def\xstart{10}
      \def\xend{45}

      % Reference line y=0
      \draw[dashed, thick] (0,0) -- (50,0);

      % Draw blue parabola:
      \draw[blue, line width=1.5pt, domain=\xstart:\xb, smooth, variable=\x, samples=200]
      plot ({\x},{\k*(\x-\xa)*(\x-\xb)});

      % Draw dashed blue parabola:
      \draw[blue, dashed, line width=1.5pt, domain=\xb:\xend, smooth, variable=\x, samples=200]
      plot ({\x},{\k*(\x-\xa)*(\x-\xb)});
      
      % Red parabola (bounced), translated in x by (x2 - x1)
      \pgfmathsetmacro{\xc}{\xb + (\xb-\xa)}      %
      \draw[red, line width=1.5pt, domain=\xb:\xend, smooth, variable=\x, samples=200]
      plot ({\x},{\k*(\x-\xb)*(\x-\xc)});
      

      % labels:
      \node[black, below] at (2, 0) {$y=0$};
      \node[black, right] at (\xend-1, 3) {$y=y_b, v_y=v_b$};
      \node[black, right] at (\xend-1,-7) {$y=y_a, v_y=v_a$};
      \node[black, below] at (\xb-2,0) {$v_y=\pm v_0$};

      \node[black, below] at (20,0) {$t \rightarrow$};
      
    \end{tikzpicture}
\caption{After a timestep $\Delta t$ shown in blue, some molecule
  trajectories pass through $y=0$, as if the ground was not there,
  ending at $y = y_a < 0$ and $v_y = v_a < 0$.  The bouncing algorithm finds $y_b$ and $v_b$, following the bounced red trajectory instead, and ending at the same time. }
\label{fig:bounce}
\end{center}
\end{figure}


There is no longer any need to bounce at the top of the cylinder, as
the last term will inevitably pull molecules back toward $y=0$.  We
will, however, have to handle particles bouncing off the bottom of the
container at $y = 0$, as illustrated in Fig.~\ref{fig:bounce}.
Suppose that after a time step $\Delta t$, we have a particle that
passed through the bottom of the container, with $v_y = v_a$ and $y =
y_a$.  We will certainly have:
$$y_a < 0$$
and
$$v_a < 0.$$ Bouncing is not so simple in this case, due to the
acceleration due to gravity.  Our approach will be to find position
$y_b$ and velocity $v_b$ if the particle had bounced, instead of
simply passing through $y=0$.  Our first step is to find the time $t$
since we passed though $y=0$, with $y = y_0 = 0$ and $v_y = v_0$.  We
do not yet know either $v_0$ or $t$, but they are easy to calculate from 1-D kinematics:
$$v_a = v_0 - gt$$
so that:
$$t = \frac{v_0 - v_a}{g}$$
and also:
$$v_a^2 = v_0^2 - 2 g y_a$$
so:
$$v_0 = - \sqrt{v_a^2 + 2 g y_a}$$
notice that we have kept the negative square root corresponding to the most
recent crossing of $y=0$ with $v_0 < 0$.  (The other root corresponds
to an earlier time with the molecule passing upwards through $y=0$.)
And so:
$$t = \frac{-v_a-\sqrt{v_a^2 + 2 g y_a}}{g}$$
That result for $t$ is completely valid, except it is numerically
unstable (...) so instead we use the equivalent form:
\begin{equation}
t = \frac{- 2 y_a}{\sqrt{v_a^2 + 2 g y_a} - v_a}
\end{equation}
which is numerically stable for $v_a < 0$.

With $v_0$ and $t$ in hand, we can calculate the position $y_b$ and the
$y$-component of veloctiy $v_b$ had the particle bounced at $y=0$, which is to say that $v_0 \to -v_0$.  So:
\begin{eqnarray*}
y_b &=& (-v_0 t) - \frac{1}{2}\; g t^2 \\
\end{eqnarray*}
and:
\begin{equation}
v_b = (-v_0 t) + gt \\
\end{equation}
Take care that $-v_0$ is positive!

Note that if $t$ is large enough, the particle may have fallen below
$y=0$, so that $y_b<0$.  You will have to keep bouncing molecules
until all of them have $y > 0$ at the end of the time step.


\begin{equation}
v_b = -v_a - 2 g (y_b - y_a)
\end{equation}

\section{Homework Problems}

{\bf Zero credit if you do not use our system of units consistently!}
When I ask you to calculate a quantity, e.g. temperature, I want the
answer in our units, e.g. $T_0$.  Do not first calculate in SI units
and then convert!\\[3pt]

\noindent
The starting point for this assigment is Problem 4 from the previous
homework.  You will want to initialize a numerical model for an ideal
gas tracking variables $v_x$, $v_y$, and $y$ for $N=10000$ molecules,
with height $H=4a$.  If you haven't already, define a function that
brings your gas into thermal contact with a reservoir of temperature
$T_R$, with an input parameter for the number of collisions.  Call it
what you want, but we will refer to it as {\tt collide\_reservoir}
below.  Use {\tt collide\_reservoir} to bring your gas into
equilibrium with a reservoir with $T_r = 1 T_0$.\\[3pt]

\noindent
{\bf Problem 0:} If you did not finish the last problem of HW 1, submit it with this assignment instead.  If you completed it already in HW 1, leave a comment to that effect in your notebook.\\[3pt]

\noindent
{\bf Problem 1:} Modify your {\tt collide\_reservoir} function to
measure the kinetic energy of the reservoir before and after the
collisions, and return the heat $Q$ transferred to the gas as the
return value for the function.  Start with the gas at $T=1~\rm T_0$.
Call {\tt collide\_reservoir} to perform $n_{\rm coll} = N = 10000$
collisions with a reservoir at $T=2~\rm T_0$.  Record the gas
temperature and the heat exchanged at this step.  Call the {\tt
  collide\_reservoir} function a total of 10 times in this manner,
recording the gas temperature and heat exchanged each time.  Plot $Q$
versus iteration (from 1 to 10) and $T$ versus iteration (from 1 to
10).  You should see the $Q$ values are large for the first few
iterations and then oscillate near zero.  You should see the
temperature rising steadily until it reaches $T=2 T_0$.  This exercise
should convince you that it takes about ten times as many collisions
as gas molecules to reliably bring your gas into thermal equilibrium.\\[3pt]

\noindent
{\bf Problem 2:} Return your gas to $T=1~\rm T_0$.  Then call {\tt
  collide\_reservoir} with $T_r = 2\;T_0$ for $n_{\rm coll} = 10 N =
100,000$ collisions, and record the temperature of your gas and the
heat exchanged.  Next, bring your gas into thermal contact with a
reservoir at $T=3\;T_0$, and record the heat exchanged.  Repeat this
process until your reach $T=5\;T_0$.  Plot the total heat exchanged as
a function of each temperature point. For example, the total heat
exchanged at temperature $T=3\;T_0$ would be the sum of the $Q$
exchanged with the reservoirs at $T_r = 2\;T_0$ and at $T_r = 3\;
T_0$. Read off the heat capacity of your gas from the plot, and
compare it to expecation in the comments.\\[3pt]

\noindent
{\bf Problem 3:} Call your current gas model ``gas A''.  Create a
second gas model ``gas B'' with the same parameters.  Intialize gas A
to $T=1 T_0$ and gas B to $T=2 T_0$.  Bring gas A and gas B into
thermal contact with each other for $n_{\rm coll} = 10 N = 100,000$
collisions.  Measure the temperature of both after the collisions and
compare with your expectations in the comments.\\[3pt]

\noindent
{\bf Problem 4:} Define a function that updates your gas for a time
step $\Delta t$ accounting for a non-zero gravity, and handling
bouncing at $y=0$, as described in the text.  Use this to turn gas A
into a model for the atmosphere with gravity and $T = 1\rm T_0$.  To
do so use\footnote{Other combinations work, but I found this to be a reliable and reasonably efficient setting.} a total of $n_{\rm coll} = 20 N = 200,000$
collisions and $n_{\rm step} = 2000$ time steps of size $\Delta t =
0.01\tau$.  During each time step, account for gravity and then
simulate thermal contact with a reservoir at $T_r = 1 T_0$ for $n_{\rm
  coll}/n_{\rm step}$ collisions.  Measure the temperature of the gas
and the total energy (kinetic and gravitational) and compare with your
expectation.\\[3pt]

\noindent
{\bf Problem 5:} Make a histogram of the $y$-position of the
molecules, and compare to the theoretical prediction in the text.\\[3pt]

\noindent
{\bf Problem 6:} Bring your atmosphere (gas A) at temperature $T=1~\rm
T_0$ into thermal contact with (gas B) at temperature $T=4~\rm T_0$.
Apply gravity to gas A but not to gas B.  Use a total of $n_{\rm coll}
= 20 N = 200,000$ collisions, spread across $n_{\rm step} = 2000$ time
steps of size $\Delta t = 0.01\tau$, as in the previous problem.
Measure the equilibrium temperature and compare to your expectation,
based on the heat capacities of the two gases.  (Make sure you do not
use {\tt collide\_reservoir} while bring the two gases into
equilibrium.)\\[3pt]

\end{document}
